\scanpage{110}
moyennant \unsure{ceux-ci|ceci}, on peut identifier donc les \(G_i\) au \(=\) groupe \(G\), les \(\unsure{Q_i|P_i}\) aux bitorseurs triviaux sous \(G,G\), \(\Gamma\) à
\[
 \Gamma_{\unsure{0|e|G}}=\{(x,y,z)\in G\deleted{\unsure{^3}}\times G\times G\mid xyz=1\}
\]
(qu'on peut résoudre en \(x,y,z\)
\[
\begin{aligned}
 x&=z^{-1}y^{-1}\\
 y&=x^{-1}z^{-1}\\
 z&=y^{-1}x^{-1}
\end{aligned}\hspace{9em}\biggr)
\]

La condition (A) signifie que \(\Gamma\) est \deleted{\unsure{bien|linéaire}}\inserted{\text{\unsure{nécess.}}} \deleted{\unsure{alors identique à}}\inserted{\text{\unsure{donné|donnée} par}}
\[
 \Gamma=\{(g_2^{\deleted{\unsure{-1}}}\deleted{\unsure{g}}g_3^{-1},g_3g_1^{-1},g_1g_2^{-1})
 \in G\times G\times G\mid g_1,g_2,g_3\in G\}
\]
Si \(\gamma\inserted{=(a_1,a_2,a_3)}\) est remplacé par
\[
\begin{aligned}
 \gamma'&=\sigma(\deleted{g}\alpha_1,\deleted{g}\alpha_2,\deleted{g}\alpha_3)(\gamma)\\
 &=(\deleted{g}\alpha_2 a_1\deleted{g}\alpha_3^{-1},\deleted{g}\alpha_3 a_2\deleted{g}\alpha_1^{-1},
 \deleted{g}\alpha_1 a_3\deleted{g}\alpha_2^{-1})
\end{aligned}
\]
alors les iso. \(\widetilde{a_{1}},\widetilde{a_{2}},\widetilde{a_{3}}\) sont remplacés par
\[
 \deleted{\operatorname{int}(\unsure{g_3|\alpha_3})^{\unsure{-1}}\widetilde{a_{1}}\operatorname{int}(\unsure{g_2|\alpha_2})}
\]
\[
 \left\{\begin{aligned}
 \widetilde a'_1&=\operatorname{int}(\alpha_3)\widetilde a_1\operatorname{int}(\alpha_2)^{-1}\\
 \widetilde a'_2&=\operatorname{int}(\alpha_1)\widetilde a_2\operatorname{int}(\alpha_3)^{-1}\\
 \widetilde a'_3&=\operatorname{int}(\alpha_2)\widetilde a_3\operatorname{int}(\alpha_1)^{-1}
 \end{aligned}\right.
\]
% archive sheet number: 111
